\begin{ExerciseBlock}
\ExerciseHeading{}{Verständnisfragen}{geometrical_and_wave_optics_checks}
Entscheiden Sie, welche Aussagen zutreffen, und begründen Sie Ihre Entscheidung.
\begin{enumerate}
\item Die geometrische Optik gilt näherungsweise, wenn die Wellenlänge viel größer als die optischen Hindernisse ist.\TrueFalseAnswer{falsch}
\item Bei gleicher Ausbreitungsgeschwindigkeit und gleicher Spaltbreite werden Wellen mit höherer Frequenz stärker gebeugt.\TrueFalseAnswer{falsch}
\item Die Maxima hinter einem Doppelspalt haben auf einem ebenen Schirm alle denselben Abstand.\TrueFalseAnswer{falsch}
\item Bei gleicher Öffnung werden Wellen mit kleiner Wellenlänge stärker gebeugt als solche mit großer Wellenlänge.\TrueFalseAnswer{falsch}
\item Totalreflexion ist nur beim Übergang vom optisch dichteren zum optisch dünneren Medium möglich.\TrueFalseAnswer{wahr}
\item Brechung und Beugung treten ausschließlich bei Lichtwellen auf.\TrueFalseAnswer{falsch}
\item Bei konstanter Absorptionsstärke nimmt die Lichtintensität im Material exponentiell mit der Weglänge ab.\TrueFalseAnswer{wahr}
\item Der Absorptionskoeffizient ist unabhängig von der Wellenlänge.\TrueFalseAnswer{falsch}
\end{enumerate}
\end{ExerciseBlock}
